\section*{第 \ref{ch:b3:facies-models} 章——相模型}

\begin{solution}{exo:b3:facies-models:1}
由海向陆依次为：离岸带、含风暴砂的离岸过渡带、\omterm{def:b3:facies-models:shoreface}{临滨}、
海滩、带沼泽的滨海平原。各环境依次覆盖在其向海一侧的
邻近环境之上，海岸线向海推进，即海退。若为海侵，则
层序以相反方向堆叠：底部为滨海平原，
顶部为离岸泥，
且常在\omterm{def:b3:facies-models:shoreface}{临滨}过渡处出现侵蚀面。
\end{solution}

\begin{solution}{exo:b3:facies-models:2}
(a) 河道充填或\omterm{def:b3:facies-models:rivers}{点沙坝}：底部为砾质冲刷面，
其上为交错层，再上为波痕和带根土的泥。
(b) 进积的\omterm{def:b3:facies-models:shoreface}{临滨}或三角洲前缘：下部为含风暴砂层与
\omterm{def:b3:facies-models:hcs}{丘状交错层理}的泥，上部为交错层理的\omterm{def:b3:facies-models:shoreface}{临滨}砂与海滩。
(c) 叠置的\omterm{def:b3:facies-models:rivers}{辫状河}砂、风成沙丘或碳酸盐浅滩：为粒度均一的交错层系，
风成情形具大型沙丘前积层，
浅滩情形则含骨骼和鲕粒。
\end{solution}

\begin{solution}{exo:b3:facies-models:3}
$\phi = (2.65 - 2.40)/(2.65 - 1.0) = 0.152$；含卤水时，
$0.25/1.55 = 0.161$。
\end{solution}

\begin{solution}{exo:b3:facies-models:4}
$e^{-2\pi z/\lambda} = e^{-\pi} = 0.043$，当 $z = \lambda/2$ 时，约为
\qty{38}{m} 和 \qty{55}{m}。在 \qty{20}{m} 处：$e^{-2\pi\times20/75} = 0.19$
和 $e^{-2\pi\times20/110} = 0.32$，均为地表运动的比例。
\end{solution}

\begin{solution}{exo:b3:facies-models:5}
$(60 - 47.6)/(189 - 47.6) = 0.088$；$(60 - 52.6)/(189 - 52.6) = 0.054$。
若骨架值的不确定度为 \qty{5}{\micro\second/ft}，孔隙度
则变化约三分之一：时间平均方程需要用岩心校验的骨架值。
\end{solution}

\begin{solution}{exo:b3:facies-models:6}
$(6.0 + 3.5 + 8.0 + 2.5)/40 = 20/40 = 0.50$。孔隙体积：$20\times0.22
\times10^6 = \qty{4.4e6}{m^3}$，每平方千米。
\end{solution}

\begin{solution}{exo:b3:facies-models:7}
靠近峡谷处仅见 Ta，上部分段已被后续水流移除；沿扇体向下，
出现 Tc（波痕）与 Td（粉砂），位于 Te（泥）之下。水流
随距离增加失去速度与粗粒载荷，仅细粒组分能搬运
\qty{30}{km} 之远，其沉积作用由较慢、形成波痕的水流完成，
而非高密度悬浮体的沉降。
\end{solution}

\begin{solution}{exo:b3:facies-models:8}
\omterm{def:b3:facies-models:rivers}{曲流河}的\omterm{def:b3:facies-models:rivers}{点沙坝}：冲刷河床底部的残留层，
深河道的交错层理砂，坝顶的波痕砂，
河漫滩的带根泥。该层序记录了从河床底到河岸的
河道，故河道深约 \qty{9}{m}（泥经压实后略浅，
）。该数值不能外推至底面以下：底面为侵蚀面，
其下可能是任何更老的沉积物。
\end{solution}

\begin{solution}{exo:b3:facies-models:9}
$60/40\,000 = 0.0015$，\ang{0.09}，为缓坡。$300/1500 = 0.20$，\ang{11.3}，
后接深 \qty{5}{m} 的平顶：为\omterm{def:b3:facies-models:ramp}{镶边陆架}，边缘有礁或浅滩颗粒，
斜坡下部则为块体与颗粒流。
\end{solution}

\begin{solution}{exo:b3:facies-models:10}
砂带经过某固定地点所需的时间为 $w/v$，在此期间，
其下方海底相对海平面以速率 $A$ 抬升：砂体充填厚度为 $h = A\,w/v$。
$2000\times0.2/20 = \qty{20}{m}$。若 $A = 0$，\omterm{def:b3:facies-models:shoreface}{临滨}仅留下
其自身剖面的厚度，形成薄层，因为各砂带彼此
叠置而不加积。比值 $A/v$ 即为轨迹角的正切，
此轨迹角见第~\ref{ch:b3:sequence-stratigraphy}~章的命题，
各相的厚度等于其宽度乘以该正切值。
\end{solution}

\begin{solution}{exo:b3:facies-models:11}
\omterm{def:b3:facies-models:wireline}{密度测井}：$(2.65 - 2.30)/1.65 = 0.212$；\omterm{def:b3:facies-models:wireline}{声波测井}：$(86 - 52.6)/136.4 = 0.245$。
取 $\phi = 0.245$，则 $\rho_{\mathrm{ma}} = (2.30 - 0.245)/(1 - 0.245) =
\qty{2.72}{g/cm^3}$，密度大于石英，这与由方解石胶结的砂岩
（方解石密度 2.71）或含方解石颗粒的砂岩相符。
\end{solution}

\begin{solution}{exo:b3:facies-models:12}
覆于煤层之上的离岸泥使不相邻的环境接触：一个
\omterm{def:b3:sequence-stratigraphy:parasequence}{海泛面}，此处为波浪侵蚀形成的\omterm{def:b3:sequence-stratigraphy:surfaces}{海侵面}，贝壳碎屑
构成残留层。随着海岸线向陆迁移，波浪将底面
削蚀至约\omterm{def:b3:facies-models:shoreface}{正常浪基面}：海滩、上\omterm{def:b3:facies-models:shoreface}{临滨}与
滨海平原顶部，约 \qty{10}{m} 厚的剖面可能
已被剥蚀缺失。
\end{solution}

\begin{solution}{pb:b3:facies-models:1}
\textbf{1.} 白垩（冷、低密度），自 \qty{766}{m} 起为页岩（热），自
\qty{866}{m} 起为石灰岩（冷、致密、快），自 \qty{1046}{m} 起为砂岩（钟形），
自 \qty{1166}{m} 起为片岩。

\textbf{2.} 白垩 \qty{766}{m}，页岩 100，石灰岩 180，砂岩
\qty{120}{m}；盖层底界位于 $441.0 - 1166 = \qty{-725}{m}$。

\textbf{3.} 近地表处为 $(2.71 - 2.03)/1.71 = 0.40$，\qty{750}{m} 处为 $(2.71 -
2.36)/1.71 = 0.20$：白垩在上覆岩层的重量下，
经压实作用损失了一半的孔隙度。

\textbf{4.} 其颗粒为密度不确定的黏土，其中的水
部分与之结合，因此该方程的假设不成立。

\textbf{5.} 结晶基底：穿过盖层后，
其下预计不存在储层。

\textbf{6.} $(2.71 - 2.573)/1.71 = 0.080$。

\textbf{7.} $(58.9 - 47.6)/141.4 = 0.080$：相同。

\textbf{8.} $(2.65 - 2.32)/1.65 = 0.200$。

\textbf{9.} $(79.9 - 52.6)/136.4 = 0.200$：相同。

\textbf{10.} $0.33/1.55 = 0.213$，多出一又三分之一个百分点。

\textbf{11.} 砂岩，其砂中孔隙度为 \qty{20}{\percent}；
石灰岩上方的页岩对其起封盖作用，泛滥平原泥
将砂岩分隔成若干块体。

\textbf{12.} 钟形：交错层理砂中自然伽马值低，向上经变细段升至生根泥的高值，
底界陡直。

\textbf{13.} \omterm{def:b3:facies-models:association}{向上变细序列}：侵蚀底界之上依次为河道底部滞留沉积、
河道沙丘（交错层理）、\omterm{def:b3:facies-models:rivers}{点沙坝}顶部（波痕）、
泛滥平原（生根泥）。

\textbf{14.} 河道、位于凹岸的\omterm{def:b3:facies-models:rivers}{点沙坝}、天然堤与
带有土壤的泛滥平原并列分布。

\textbf{15.} 约为砂层与变细段厚度之和，$10.4 + 3.2 =
\qty{13.6}{m}$，压实前略厚一些。

\textbf{16.} \omterm{def:b3:facies-models:rivers}{曲流河}：盆地沉降期间，
河道带因\omterm{def:b3:facies-models:rivers}{河道改道}而反复迁移，
每条新河道都切入前一条河道遗留的泛滥平原。

\textbf{17.} 不能：每个底界都是侵蚀面，其下的泛滥平原泥与其上的河道砂，
未必曾经相邻。

\textbf{18.} $(60 - 30)/(105 - 30) = 0.40$。

\textbf{19.} 最薄情形：$14\times(0.60 + 0.40\times0.25) = \qty{9.80}{m}$；
最厚情形：$16\times(0.65 + 0.40\times0.20) = \qty{11.68}{m}$。

\textbf{20.} 纯砂 \qty{58.8}{m}；变细段 \qty{33.4}{m}；泥
\qty{27.8}{m}。

\textbf{21.} 净砂厚度 $= 58.8 + 0.40\times33.4 = \qty{72.2}{m}$；\omterm{def:b3:facies-models:ntg}{净毛比}
$72.2/120 = 0.60$。

\textbf{22.} 比例取 0.20 与 0.60 时，厚度分别为 \qty{65.5}{m} 和 \qty{78.8}{m}，
仅因截断值的选取，二者相差 \qty{13}{m}。

\textbf{23.} $72.2\times0.20\times10^6 = \qty{1.44e7}{m^3}$，
每平方千米。

\textbf{24.} 瓦尔登-2 井中三叠纪砂岩的净砂厚度为
\textbf{\qty[text-series-to-math=true]{72.2}{m}}。
\end{solution}
